\section{Execution-Based Results}

Execution-based evaluation was performed using CodeEditorBench, where the
generated C++ translations were compiled and executed against predefined test
cases. Compared to BLEU score, this evaluation provides a more direct measure of
functional correctness because it captures whether the translated program
compiles and produces the expected outputs.

\subsection{Accepted Solutions and Compilation Errors}

\begin{table}[htbp]
\centering
\caption{Execution-based evaluation results}
\label{tab:execution_results}
\resizebox{\textwidth}{!}{%
\begin{tabular}{lcccc}
\hline
Model & Accepted & Wrong Answer & Runtime Error & Compilation Error \\
\hline
StarCoder2 Base & 55 & 57 & 4 & 390 \\
StarCoder2 Fine-Tuned & 133 & 75 & 13 & 287 \\
Qwen Base & 101 & 63 & 15 & 324 \\
Qwen Fine-Tuned & 168 & 96 & 9 & 219 \\
\hline
\end{tabular}%
}
\end{table}

The number of accepted solutions increased significantly after fine-tuning for
both model families. %szb: ugyanaz, amit az előbb írtam

For StarCoder2, the number of accepted translations increased from 55 for the
base model to 133 for the fine-tuned model. This corresponds to an improvement
of approximately 142\%. At the same time, the number of compilation errors
decreased from 390 to 287.

%szb: ezekből a részekből (egész results) abszolút hivatkoznak a hivatkozások, vagyis olyanok hogy Table X shows the results of ...

The Qwen models achieved stronger performance overall. The base Qwen model
produced 101 accepted solutions, while the fine-tuned version reached 168
accepted solutions. This represents an improvement of approximately 66.3\%.
Compilation errors also decreased substantially, from 324 to 219.

Compilation error remained the most frequent failure type for all evaluated
models. Even after fine-tuning, compilation errors represented the largest
portion of failed translations. 

%szb: fura, mert efelett már írsz az error type-kről, majd a kövi subsection-ban is. Ha már mindenképp ketté akarjuk bontani, akk a következőképpen: 
%- puszta eredmények, melyik modell hogy teljesített, minden error-t egyként kezelve
% - rontások különböző fajtái (error types)






\subsection{Other Error Types}

\begin{table}[htbp]
\centering
\caption{Additional execution error types}
\label{tab:error_types}
\resizebox{\textwidth}{!}{%
\begin{tabular}{lccc}
\hline
Model & Presentation Error & Time Limit Exceeded & Output Limit Exceeded \\
\hline
StarCoder2 Base & 2 & 2 & 0 \\
StarCoder2 Fine-Tuned & 0 & 2 & 0 \\
Qwen Base & 3 & 4 & 0 \\
Qwen Fine-Tuned & 10 & 5 & 3 \\
\hline
\end{tabular}%
}
\end{table}

In addition to accepted solutions and compilation errors, several other failure
categories appeared during evaluation.

Wrong Answer (OJ\_WA) was the second most common error type. For StarCoder2,
the number of wrong answers increased from 57 in the base model to 75 after
fine-tuning. Similarly, Qwen increased from 63 to 96 wrong answers after
fine-tuning.

Runtime Error (OJ\_RE) occurred relatively infrequently. For StarCoder2, runtime
errors increased from 4 to 13 after fine-tuning, while for Qwen they decreased
from 15 to 9. 

Time Limit Exceeded (OJ\_TL) errors were rare for all models, indicating that
most generated solutions were computationally efficient enough to finish within
the allowed execution time. Presentation errors and output limit errors also
occurred only occasionally and had limited influence on the overall results.


